\section{Localization of arbitrary approximating domains}
\label{loc:section}
\label{sec:localization}

Marked kernel convergence to a bounded domain does not bound the
diameters, areas, or numbers of remote appendages of its
approximants. Thus Proposition~\ref{uni:bounded} does not yet prove
the full theorem. We reduce the general case to that proposition
by a finite edge exploration. The exploration will leave an
ordinary bounded Dobrushin problem, with probability tending to
one, and change the transferred curve by a vanishing amount.
Every probability estimate below uses rectangles in a fixed
bounded physical region. The deterministic gates belong to the original
domain and organize the probability estimates; their shadows are the
components separated from the basepoint. After the endpoint heads have
been explored, portals bound monochromatic caps, which a cluster search
shields with high probability while preserving the conditional FK law
on success.

Let \((\Omega_n,a_n,b_n)\) be an arbitrary sequence in the theorem,
with mesh \(\delta_n\to0\). Translate the fixed kernel
basepoint to \(0\), and write
\[
 \chi_n:\Omega_n\longrightarrow\D,\qquad
 f_n=\chi_n^{-1}
\]
for the normalized disc maps. Their kernel limit is
\(\chi:\Omega\to\D\), with inverse \(f\). Fix once and for all
\[
 S>\sup_{z\in\Omega}|z|+10.
\]
All subsequent subsequences retain this choice of \(S\).

\subsection{Two uniformization lemmas}

\begin{lemma}[Small parameter shadows of remote pieces]
\label{loc:remote}
For every connected interior arc \(E_n\subset\Omega_n\setminus
B(0,S-1)\), its disc image has diameter tending to zero,
uniformly over the choice of arc. The same is true for the
noncentral shadow of a crosscut contained there.

More precisely, fix a point \(\xi\in\partial\D\), a positive
parameter tolerance \(a\), and a positive physical tolerance
\(\rho\). For all sufficiently large \(n\), one can choose
\(s_n\) in a fixed interval \([s_-,s_+]\subset(0,a)\) such that
the entire preimage of
\(\partial B(\xi,s_n)\cap\D\) is contained in \(B(0,S-2)\)
and has length at most \(\rho\).
The endpoints and shadows are understood in the prime-end
compactification. The conclusion holds simultaneously for
finitely many specified \(\xi\)'s, with disjoint parameter caps
when those points are distinct.
\end{lemma}

\begin{proof}
Choose as many consecutive fixed dyadic rings about \(\xi\)
as needed, all with radii below \(a\). For a radius \(s\) in
their range, the circular crosscut contains the inward radial
point \((1-s)\xi\). All these radial points belong to a fixed
compact subset of \(\D\), because the smallest radius is
positive. Compact convergence of \(f_n\) puts their physical
images in a fixed neighborhood of the bounded limit domain.
In particular they are separated from \(\partial B(0,S-2)\)
by a positive gap \(g\), uniformly for all large \(n\).

Let \(L_n(s)\) be the length of only that part of
\(f_n(\partial B(\xi,s)\cap\D)\) lying in \(B(0,S-2)\).
Polar Cauchy--Schwarz and univalence imply
\begin{equation}
 \int L_n(s)^2\,\frac{\dd s}{s}
 \le 2\pi
 \int_{f_n^{-1}(B(0,S-2))}|f_n'(z)|^2\,\dd A(z)
 \le 2\pi^2(S-2)^2 .
 \label{loc:restricted-area}
\end{equation}
The right side bounds only the area of a fixed physical ball.
It makes no assumption on \(\area(\Omega_n)\).
By taking enough fixed dyadic rings, a regular radius can be
selected for which \(L_n(s)<\min(\rho,g)\).
If its full preimage arc left the physical ball, the segment
from the anchored radial point to its first exit would already
have length at least \(g\) inside the ball, contradicting this
inequality. The whole arc therefore stays in the ball and has
length at most \(\rho\). Radii can be chosen away from the
finitely many contour degeneracies of the polygon.

For a fixed prescribed parameter accuracy, apply this
construction to a finite collection of circle points whose
small caps cover \(\partial\D\). Images of remote arcs lie
uniformly close to \(\partial\D\): every fixed compact subset
of \(\D\) has bounded physical preimage contained in
\(B(0,S-1)\) for all large \(n\). They avoid the bounding
crosscuts of the chosen caps, since those have physical
preimages in \(B(0,S-2)\). Connectedness therefore traps each
remote image in one of the small caps. This proves the
uniform diameter assertion. If the remote arc is a crosscut,
join its endpoints by the short boundary-circle interval.
Its noncentral shadow is contained in the same cap.

For finitely many distinct prescribed points, first choose
disjoint outer parameter neighborhoods and then perform the
selection within each. Every radius range remains fixed at
this stage, although its lower endpoint can be very small.
All mesh limits are taken after these ranges have been fixed.
\end{proof}

\begin{lemma}[Small attached changes preserve normalized coordinates]
\label{loc:obstructions}
Let \(V_n\subset\D\) be the simply connected component containing
\(0\) after removing finitely many closed, piecewise analytic
obstructions and filling the other components. Assume that
the obstructions lie within distance \(\varepsilon_n\to0\)
of \(\partial\D\), and that every point on an obstruction
can be joined to \(\C\setminus\D\) within the removed set by
a connected set of diameter at most \(\varepsilon_n\).
Then the normalized maps \(F_n:\D\to V_n\) satisfy
\begin{equation}
 \sup_{\overline\D}|F_n(z)-z|\longrightarrow0.
 \label{loc:maps-identity}
\end{equation}
For every prime-end occurrence \(w=F_n(\zeta)\), this also
gives \(|F_n^{-1}(w)-w|\to0\), uniformly. The same conclusion
holds for any number of such obstructions; no bound on their
number or total length is required.
\end{lemma}

\begin{proof}
First note that filling preserves a small attachment property.
The central component contains \((1-\varepsilon_n)\D\).
A point in a discarded component lies in the remaining
boundary collar. Follow its outward radial segment to its
first encounter with an obstruction or the circle. Before
that encounter the segment stays in the same discarded
component and has length at most \(\varepsilon_n\).
Append the obstruction's small attachment when necessary.
Every point of the filled complement thus has a connected
route of diameter at most \(2\varepsilon_n\) to the old
outer complement.

We give a quantitative crosscut proof of
\eqref{loc:maps-identity}. Suppress \(n\) and suppose
\[
 (1-d)\D\subset V\subset\D,\qquad
 \diam A_x\le e
\]
for a connected complement attachment \(A_x\) from each
complement point \(x\) to the old exterior. The normalized
map \(F:\D\to V\) extends continuously to the physical
boundary, since \(V\) has a finite, locally connected
piecewise analytic boundary, with occurrences retained at
slits. Put \(a=F'(0)\). Schwarz's lemma, also applied to
\(F^{-1}\) on \((1-d)\D\), gives \(1-d\le a\le1\).
Schwarz--Pick applied to \(F(z)/z\) yields
\begin{equation}
 |F(z)-z|\le\frac{2d}{1-|z|}.
 \label{loc:interior-map}
\end{equation}

Fix \(\zeta\in\partial\D\) and \(0<h<1\).
Length--area on the circular crosscuts
\(\partial B(\zeta,t)\cap\D\), \(h<t<\sqrt h\),
gives one whose image \(\sigma\) has length at most
\[
 L_h=\frac{2\pi}{\sqrt{\log(1/h)}}.
\]
It contains \(F((1-t)\zeta)\), whose distance from \(\zeta\)
is at most \(\sqrt h+2d/h\) by
\eqref{loc:interior-map}. Let \(p,q\) be the endpoint
occurrences of \(\sigma\). Their physical locations have
distance at most \(L_h\). Join them in the complement by
their \(e\)-attachments and the shorter arc of
\(\partial\D\) between the attachment endpoints.
This gives a connected complement set within distance
\(L_h+3e\) of \(p\).

The complement is locally path connected in the present
finite-arc setting. For any \(\epsilon>0\), the joining
continuum can therefore be replaced by a simple complement
arc within its \(\epsilon\)-neighborhood. Together with
\(\sigma\) it forms a Jordan curve. If the physical endpoints
coincide, \(\sigma\) itself is the corresponding closed curve;
the same separation argument applies to its two boundary
occurrences. When
\[
 2L_h+\sqrt h+2d/h+3e+\epsilon<1,
\]
this curve is in a ball missing \(0\). Its bounded side
contains the noncentral image cap. It follows that
\[
 |F(\zeta)-\zeta|
 \le2L_h+\sqrt h+2d/h+3e+\epsilon .
\]
Let \(\epsilon\downarrow0\). The bound is uniform in
\(\zeta\), and the maximum principle gives the same bound
on \(\overline\D\). Choose \(h\) small first and then
\(d,e\) small. This proves \eqref{loc:maps-identity}.
The argument uses the number of obstructions nowhere.

Finally, if \(w=F_n(\zeta)\), then
\(|F_n^{-1}(w)-w|=|\zeta-F_n(\zeta)|\).
This proves the inverse statement with prime-end
occurrences, without making an inverse single-valued at
an unlabelled physical slit point.
\end{proof}

\begin{remark}
\label{loc:collars}
Both lemmas apply to the cell collars used below.
In the initial disc, each retained cell sector has vanishing
diameter by Lemma~\ref{geo:crosscut}. A boundary cell is
attached to the old boundary through its own sector.
Filling behind a boundary pinch is covered by the first
paragraph of Lemma~\ref{loc:obstructions}.
The finite-arc boundary condition in that lemma is imposed
only on the approximating polygons and their conformal
images, not on the limiting domain \(\Omega\).
\end{remark}

\subsection{Mandatory gates in a bounded physical band}

Set
\begin{equation}
 s_j=S+1-\frac1{j+1},\qquad
 1\le j\le N_n,\qquad
 N_n=\left\lfloor\log_2(1/\delta_n)\right\rfloor .
 \label{loc:levels}
\end{equation}
Perturb the levels, if necessary, by a fixed small fraction
of their gaps so that all circular intersections are
transverse to graph edges and avoid vertices. The estimates
below are unchanged.

\begin{lemma}[Visible gates and their ancestors]
\label{loc:gates}
Let \(U_{n,j}\) be the component of
\(\Omega_n\cap B(0,s_j)\) containing \(0\).
Its visible exit arcs on \(\partial B(0,s_j)\) are crosscuts
of the initial domain, with pairwise disjoint noncentral
shadows. Every point of \(\Omega_n\) outside that circle
lies behind one mandatory visible gate: every path from
\(0\) to that point crosses it. A gate at a later level has
one specified ancestor gate at every earlier level.

Call a gate at level \(j\) large when its length is at least
\(2^{-j}\). There are at most \(K_S2^j\) large gates at that
level. A chain of small ancestors at levels \(j\) gives
disjoint annuli centered on those ancestors, with inner
radii \(K2^{-j}\) and outer radii \(c_0(j+2)^{-2}\),
for all sufficiently large \(j\). Every path from a
circle-\(S\) portal to a point of radius \(S+2\), continued
to \(0\) within \(B(0,S)\), crosses these annuli.
\end{lemma}

\begin{proof}
For all large \(n\), none of the indicated circles is
wholly contained in \(\Omega_n\). Otherwise simple
connectivity would put its full interior in \(\Omega_n\),
contradicting kernel convergence to the much smaller
bounded domain. At a generic circle, its components
inside a finite polygon are open arcs with endpoints
on boundary occurrences. Those incident to \(U_{n,j}\)
are its visible gates.

Cut the domain at all these circular arcs. The adjacency
graph of the resulting components, with one edge for each
crosscut, is a finite tree: successively adding disjoint
proper crosscuts splits one simply connected component
into two and preserves connectedness and absence of cycles.
Root this tree at \(U_{n,j}\). The visible gates are exactly
its root-adjacent edges, and their noncentral shadows are
the corresponding disjoint branches. Every point outside
the circle belongs to a nonroot component, so the unique
first edge of its root path is its mandatory gate.
Crosscut separation makes every domain path to that point
cross this gate. The tree uses boundary occurrences, so
the argument also applies at repeated physical vertices.

A connected later gate lies outside every earlier circle
and hence in one unique earlier shadow. For two disjoint
crosscuts the root-oriented shadows are either disjoint or
nested, by the same tree argument. These ancestors are
therefore consistent across levels. A shadow may contain
further components inside the earlier circle; this does not
affect the root path or the mandatory-gate property.

The total length of gates on a circle is at most
\(2\pi(S+1)\), proving the count for large gates.
A small gate fits in a ball of radius \(K2^{-j}\)
centered at any one of its points. Successive level
gaps are comparable to \((j+2)^{-2}\). Choose \(c_0\)
sufficiently small that the outer balls lie in
disjoint radial bands, and omit finitely many initial
levels so that their inner radii are smaller than the
outer radii.

The center-to-far-point route must cross each mandatory
gate. Its appended core portion has radius at most \(S\),
whereas all gates have larger radius. Thus each crossing
is on the original portal-to-far-point path. Its two
endpoints lie outside every indicated outer ball, so
visiting the small gate requires an annular crossing.
At lattice scale the path and gate may be moved within
incident cells. Since \(2^{-j}\ge\delta_n\) for
\(j\le N_n\), these changes are absorbed by the fixed
constant in the inner radius. The smallest radial gap
between outer bands is at least \(c(N_n+2)^{-2}\).
Since \(\delta_n(N_n+2)^2\to0\), any fixed number of
buffer cells preserves their disjointness uniformly
through the moving final level \(N_n\).
\end{proof}

\subsection{Preparing monochromatic portals}

Fix an error tolerance for this stage of the argument and a finite
integer \(J\). At the levels \(j\le J\), there are only finitely
many large gates, with a bound depending on \(S,J\) but not on
\(n\). By Lemma~\ref{loc:remote}, each such gate and its shadow
has vanishing disc diameter. Pass to a further subsequence,
labeling the finitely many gates after stabilizing their numbers,
so that their locations converge on \(\partial\D\).
Let \(H\) be the resulting finite set of limiting locations;
coincident limits are recorded only once.

We shall use two kinds of cuts: the physical circle
\(\partial B(0,S)\), and small parameter caps about \(H\).
Their preimages will be disjoint, because the latter cuts can be
chosen inside \(B(0,S-2)\). The small caps are chosen to contain
all large gates at the first \(J\) levels. Before fixing their
sizes, one must remove the possibility that a current mark lies
on their distal side.

\begin{lemma}[Protected endpoint explorations]
\label{loc:heads}
For every sufficiently small parameter tolerance \(a>0\),
the two ends can be explored, each stopped before its first
unpeeled cell on a protected fence, with the following
properties:
\begin{enumerate}
\item Each explored head has disc diameter at most \(Ka+o(1)\).
The two heads lie in disjoint endpoint neighborhoods.
\item The center has unqueried connected access chains to the
two resulting front corners. These chains lie in
\(B(0,S-1)\) and on the inner side of every proposed
small \(H\)-cap and of the circle at \(S\).
\item Conditional on these local revelations, the connecting
problem is an ordinary Dobrushin problem. In the initial
disc its changes are attached obstructions of diameter
at most \(Ka+o(1)\), with fillings behind them.
\end{enumerate}
The preimages of the \(H\)-cap cuts can subsequently be
chosen to have any prescribed small physical diameter
\(\rho>0\), while retaining these properties and containing
the large gates at levels \(j\le J\).
\end{lemma}

\begin{proof}
The marked parameter endpoints converge to two distinct circle
points. Around each choose a small circular cap of angular
size at most \(a\). These two caps are disjoint. Choose their
bounding circular crosscuts using Lemma~\ref{loc:remote},
so that their physical preimages lie in \(B(0,S-2)\).
Call them endpoint fences. Their parameter radii are in
fixed intervals and can be chosen to avoid every point of
\(H\) on the fences themselves. A point of \(H\) coinciding
with an endpoint is inside its fence. Other points of
\(H\) can be kept away by reducing the endpoint neighborhood
first.

Join a middle point of each fence to \(0\) by fixed compact
paths in the parameter disc. Their physical preimages stay
in \(B(0,S-2)\) for all large \(n\), by compact convergence.
The paths, together with the fences, form two protected
access chains. Now choose the still smaller disjoint
caps about the points of \(H\), applying
Lemma~\ref{loc:remote} within fixed radius intervals.
Their physical cuts have diameter at most \(\rho\).
They avoid the protected paths and fences, since those
are compact and their boundary endpoints are away from
\(H\). When an \(H\)-point is an endpoint limit, its cap
is placed entirely behind that endpoint's fence.
For all large \(n\), the finitely many large gates lie
inside their designated \(H\)-caps.

Replace each protected fence and access path by the chain
of cells it crosses, including terminal boundary
half-cells. A generic small perturbation avoids vertices,
so the chain is connected through sides. Its disc
distance from the original path tends to zero by
Lemma~\ref{geo:crosscut}. Explore from an endpoint until
just before taking the turn in the first cell of its
protected chain, even if that turn would be forced.
No tile of that chain has been clipped.

The stopping event is unavoidable for a path to the
opposite endpoint, since that path must cross the fence.
Before the stopping event, the explored head stays on
the endpoint side of the fence, up to its neighboring
cells. Thus its parameter diameter is at most \(Ka+o(1)\).
No discarded component can contain the whole protected
chain: its connected access has not been clipped, and
the continuing path must still reach the fence. In
particular the center-to-front-corner access remains
on the continuing side. The two explorations are in
disjoint endpoint neighborhoods and may be carried out
in either order.

Lemma~\ref{geo:domain-markov} now gives an ordinary
Dobrushin problem with the two front corners as its
change points. The procedure has stopped on the first
local arrival; it has not conditioned the unknown
connecting path to hit an interior target.
All head changes are within the two small parameter
caps and are attached to the old boundary there.
The active-side convention fills only components
behind these obstructions. Their attachments and
fillings satisfy Lemma~\ref{loc:obstructions}.
\end{proof}

In the post-head domain, intersect the physical circle at \(S\)
and the preimages of the \(H\)-cap cuts with the domain.
Take the crosscut components visible from the central component
containing the two protected accesses. Call them \emph{portals}.
They are disjoint; a single cutting line may contribute many
portal fragments. Their distal caps have disjoint shadows and
separate the central component from every route to
\(|z|>S+1\). Both moving marks are on the central side.
Consequently the distal boundary interval of every portal
contains no boundary-change point and has a single color.
This prime-end assertion, rather than a statement about
unlabelled physical vertices, is what will permit exact FK
shielding.

\subsection{A finite shielding exploration}

Figure~\ref{loc:fig-shield} illustrates the retained region and the
exposed frontier in the shielding construction.

\begin{lemma}[Exact shielding of a monochromatic cap]
\label{loc:shield}
Let a distal cap have boundary color \(c\), with its portal
strictly inside the physical cutoff \(S+3\).
There is a finite edge exploration on its distal side with
the following alternatives.
\begin{enumerate}
\item On success, all retained parts of the cap have radius
less than \(S+3+K\delta_n\). The new retained boundary is
made of \(c\)-open diagonals attached to the original
\(c\)-bank. Every queried edge lies outside or on that
boundary, so the retained connecting problem has its
ordinary Dobrushin conditional law.
\item Failure implies an actual opposite-color path from
radius at least \(S+3-K\delta_n\) to cells incident to
the portal, with an endpoint error of at most
\(K\delta_n\).
\end{enumerate}
The same statements hold for all caps explored in a fixed
order, stopping the entire surgery at its first failure.
Conditioning on the complete successful revelation record
imposes no further event on the unqueried central edges.
\end{lemma}

\begin{proof}
Declare all opposite-color vertices incident to cells of
the cap reaching \(|z|\ge S+3\) to be seeds, and reveal
the exterior cells. The seed region is topologically
attached to the original cap bank. Indeed, from any
point of its exterior component, an outward radial
route reaches its first exit from the finite polygon
while staying outside the cutoff; it cannot encounter
the portal, which lies strictly further in. Include
the cells of such routes in the exterior region when
recording its boundary. This is a topological
attachment of an exposed set. The seeds need not
belong to one open opposite-color cluster.

Search every opposite-color cluster meeting this
region, revealing incident edges until its full
inner frontier is known. Stop with failure if the
search or its necessary clipping cells reach the
portal buffer. On success, every edge at the inner
frontier is closed for the searched color, hence is
an open diagonal of color \(c\), or is an original
\(c\)-bank edge. The union of the seed region and
the searched clusters is attached to the bank.
Take the component containing the portal on its
unsearched side, and fill the holes discarded on
the other side. Its new boundary is therefore a
collection of \(c\)-paths attached to that same
bank, not a collection of interior wired islands.

Every tested edge is either in the removed
seed-and-cluster region or is one of the known
frontier edges. None is strictly in the retained
unqueried component. Conditioning on the full
record, the FK spatial Markov property imposes
exactly its original Dobrushin wires, extended
by the new \(c\)-frontier. The retained interface
cannot cross that frontier into the discarded
side. The event of success is already determined
by the record, so it imposes no extra condition
on retained edges.

For a failure witness, choose a reached vertex in
the portal buffer and trace its searched-color
cluster back to its first exterior seed. That seed
is at radius at least \(S+3-K\delta_n\). The path
between it and the last buffered cells consists
of actual open edges; no artificial connection
within the whole seed region is needed. If a
last incident edge triggers failure before being
queried, the known path ends only \(K\delta_n\)
short of the portal, which gives the stated
endpoint allowance.

All constructions can be performed on incident
sectors before identifying repeated boundary
vertices. Frontiers of a multiply armed cluster
or of a fragmented portal are treated on their
accessible side. Filling holes removes them
from the retained domain and does not query
inside it. Apply the construction to caps in
a deterministic order. On success in an
earlier cap only a wired frontier outside
the remaining core has been added, so the
same argument restarts in every later cap.
\end{proof}

\begin{figure}[!htb]
\centering
\begin{tikzpicture}[x=0.82cm,y=0.82cm,>=Stealth]
 \fill[blue!8] (6.6,0) rectangle (8,3);
 \draw[very thick,red!65!black] (1.6,0)--(8,0)--(8,3)--(1.6,3);
 \draw[densely dashed,thick] (1.6,0)--(1.6,3);
 \draw[densely dashed,gray] (6.6,-0.25)--(6.6,3.25);
 \draw[very thick,red!65!black]
 (4.7,0)--(4.9,0.4)--(4.55,0.7)--(4.8,1.0)--(4.45,1.35)
 --(4.65,1.75)--(4.5,2.1)--(4.9,2.45)--(4.7,3);
 \draw[thick,blue!65!black] (7.5,0.6)--(6.8,0.7)--(6.1,0.5)--(5.3,0.8);
 \draw[thick,blue!65!black] (7.4,1.5)--(6.7,1.6)--(6.0,1.4)--(5.1,1.6);
 \draw[thick,blue!65!black] (6.0,1.4)--(5.5,1.1);
 \draw[thick,blue!65!black] (7.4,2.5)--(6.8,2.3)--(6.0,2.6)--(5.3,2.3);
 \node[align=center] at (0.15,1.5) {protected\\core};
 \draw[->] (0.9,1.5)--(1.45,1.5);
 \node[below] at (1.6,-0.15) {portal};
 \node[align=center] at (3.0,1.5) {unqueried\\retained side};
 \node[above,red!65!black] at (4.7,3.3) {new \(c\)-frontier};
 \node[above] at (7.2,3.75) {exterior seeds};
 \node[below] at (6.6,-0.15) {cutoff \(S+3\)};
 \node[below,red!65!black] at (4.0,-0.7) {original bank has color \(c\)};
\end{tikzpicture}
\caption{Schematic shielding of one monochromatic cap.
The searched clusters, drawn in blue on the distal side, grow
from the exposed exterior region.
Their tested inner frontier is open in the bank color \(c\).
On success, all queried edges are on the frontier or to its
distal side; the retained region keeps an ordinary FK law.
The picture illustrates the separation argument and is not
an additional geometric assumption.}
\label{loc:fig-shield}
\end{figure}
\FloatBarrier

\Needspace{14\baselineskip}
\subsection{A uniform bound on the failure probability}
\begin{lemma}[The gate product estimate]
\label{loc:failure}
With the first \(J\) levels and their finite set \(H\) fixed,
let all \(H\)-cut preimages have diameter at most \(\rho\).
Uniformly conditional on the two explored heads, the probability
that the shielding procedure fails is bounded by
\begin{equation}
 \begin{aligned}
 &K|H|\rho^\alpha\\
 &\quad+K\sum_{k>J}2^k
     \prod_{j_*\le j<k}K\bigl(2^{-j}(j+2)^2\bigr)^\alpha\\
 &\quad+K\delta_n^{-2}
     \prod_{j_*\le j\le N_n}K\bigl(2^{-j}(j+2)^2\bigr)^\alpha
 +o(1).
 \end{aligned}
 \label{loc:failure-bound}
\end{equation}
Here \(j_*\) and the constants are independent of \(J,n\)
and of the unbounded part of the domain. In particular,
the failure probability tends to zero by first taking
\(J\) large, then \(\rho\) small at that stage, and finally
the mesh sufficiently small.
\end{lemma}

\begin{proof}
We compare all possible failure arms of one portal type
and one bank color together. Inside an annulus separated
from the portals and from the two ends of an arm, their
accessible side banks are free for the intruding color.
Lemma~\ref{geo:arm} applies to this union, including all
portal fragments and boundary-sector occurrences.
The comparison wires radial ends only; it counts actual
open arms and permits arbitrary off-shell conditioning.
It therefore gives product bounds in disjoint buffered
annuli without a factor for the number of fragments.
Earlier successful searches merely add favorable
single-color frontiers and do not change this argument.

If failure ends at a portal on an \(H\)-cut, the actual
intrusion arm from Lemma~\ref{loc:shield} joins a ball
containing that entire cut to a fixed outer radius,
say \(1/2\). The source is outside \(S+3-O(\delta_n)\),
whereas the cut lies inside \(S-2\). Applying
Lemma~\ref{geo:arm} to all fragments of that cut and
summing over the finitely many \(H\)-cuts gives the
first term in \eqref{loc:failure-bound}.

Now suppose failure ends at a circle-\(S\) portal.
Join its endpoint to \(0\) through the central core,
which excludes the initial \(H\)-caps and stays
inside radius \(S\), up to the fixed cell buffer.
Select a far point on the arm at radius \(S+2\).
This center-to-far-point route has the mandatory
initial gate sequence of Lemma~\ref{loc:gates}.

If any gate at a level \(j\le J\) is large, it lies
inside its selected parameter \(H\)-cap. The arm
therefore visits that cap, while its core extension
stays on the inner side. It crosses the corresponding
\(H\)-cut. This gives an arm through an annulus about
the physically tiny cut, again up to a bounded number
of cells. Circle-\(S\) portals are outside these
annuli because the \(H\)-cut lies inside \(S-2\).
The union comparison for the circle-portal type
therefore gives the same first term.

Otherwise let \(k>J\) be the first level at which
the gate is large. There are at most \(K2^k\)
choices for it, and each fixes all its ancestors.
The earlier ancestors are small.
Figure~\ref{loc:fig-gates} illustrates their surrounding annuli. For
\(j_*\le j<k\), the same arm crosses an annulus
with inner scale \(K2^{-j}\) and outer scale
\(c_0(j+2)^{-2}\). These annuli lie in disjoint
radial bands and miss the circle-\(S\) portals.
Both endpoints of the arm are outside the bands.
Applying the conditional product bound gives
\[
 K2^k\prod_{j_*\le j<k}
       K\bigl(2^{-j}(j+2)^2\bigr)^\alpha .
 \]
The two possible bank colors cost only another
fixed factor.

Finally, if all gates through \(N_n\) are small,
specify the connected initial-domain sector of an incident
cell at the first crossing of radius \(S+2\).
That sector fixes the ancestor chain, since it is connected
and lies outside all the tested circles. The angle bound
gives cell areas at least \(c_\eta\delta_n^2\)
and bounded local degree. A bounded annular band
therefore has at most \(K\delta_n^{-2}\)
possible terminal cells and sectors. The product
through \(N_n\) gives the last term.

For completeness, the logarithm of the product
through \(k\) is
\[
 -\frac{\alpha\log2}{2}k^2+O(k\log(k+2)).
 \]
It dominates both the factor \(2^k\) and,
when \(k=N_n\), the factor
\(\delta_n^{-2}=\exp(O(N_n))\).
The last term vanishes, and the tail sum can be
made arbitrarily small by increasing \(J\).
At every fixed stage the number \(|H|\) is finite,
so \(\rho\) can then make the first term small.
Cell buffers and generic contour perturbations
are absorbed once the stage is fixed and the
mesh is sufficiently small. All comparison
annuli have uniformly bounded physical radii.
\end{proof}

\begin{figure}[tb]
\centering
\begin{tikzpicture}[x=0.9cm,y=0.9cm,>=Stealth]
 \fill[gray!6] (1.0,-1.4) rectangle (2.6,1.4);
 \fill[gray!6] (3.7,-1.4) rectangle (5.1,1.4);
 \fill[gray!6] (6.1,-1.4) rectangle (7.3,1.4);
 \draw[->] (0,-1.75)--(8.2,-1.75) node[right] {outward};
 \foreach \x/\lab in {1.8/{s_j},4.4/{s_{j+1}},6.7/{s_{j+2}}}
 {
   \draw[densely dashed,gray] (\x,-1.4)--(\x,1.4);
   \node[below] at (\x,-1.75) {\(\lab\)};
 }
 \draw[blue!65!black,thick,->]
 (0,0.1)..controls(0.8,0.6)and(1.2,0.2)..(1.8,0.1)
 ..controls(2.8,-0.7)and(3.6,-0.5)..(4.4,-0.2)
 ..controls(5.1,0.6)and(6.1,0.5)..(6.7,0.35)
 ..controls(7.2,-0.1)and(7.7,0.2)..(8,0.25);
 \draw[red!65!black,very thick] (1.8,-0.09)--(1.8,0.29);
 \draw[red!65!black,very thick] (4.4,-0.34)--(4.4,-0.06);
 \draw[red!65!black,very thick] (6.7,0.25)--(6.7,0.45);
 \draw (1.8,0.1) circle (0.7);
 \draw (1.8,0.1) circle (0.23);
 \draw (4.4,-0.2) circle (0.6);
 \draw (4.4,-0.2) circle (0.17);
 \draw (6.7,0.35) circle (0.5);
 \draw (6.7,0.35) circle (0.12);
 \node[above] at (1.8,1.5) {small ancestor};
 \node[above] at (5.5,1.5) {disjoint buffered bands};
 \node[align=center] at (4.4,-2.6)
 {inner scale \(2^{-j}\), outer scale \(c_0(j+2)^{-2}\)};
\end{tikzpicture}
\caption{Schematic gate hierarchy, with the radial direction
flattened for clarity. Short segments mark gates, pairs of circles
bound their annuli, and shaded strips mark disjoint radial bands.
The curved intrusion arm visits each mandatory
small ancestor and therefore crosses every surrounding annulus.
The exponentially small inner scales and polynomially spaced
outer scales give the product in \eqref{loc:failure-bound}.}
\label{loc:fig-gates}
\end{figure}
\FloatBarrier

\subsection{Reduction, mixtures, and the full theorem}

\begin{proposition}[Reduction to bounded Dobrushin problems]
\label{loc:reduction}
From every admissible approximation sequence one can extract
a further subsequence and perform finite edge revelations
with success probability tending to one, so that on success:
\begin{enumerate}
\item the unqueried connecting domain \(\widehat\Omega_n\)
is a Dobrushin polygon on the original mesh with the
same thermal parameters and diameter at most \(2S+20\);
\item conditional on the complete revelation record, its
interface has the ordinary Dobrushin FK law;
\item every deterministic sequence of possible successful
outputs \((\widehat\Omega_n,\widehat a_n,\widehat b_n)\)
converges in the marked Carath\'eodory sense to
\((\Omega,a,b)\);
\item the normalized coordinates on the retained domain
are uniformly \(o(1)\) from its original normalized
coordinates, including prime-end occurrences, and
the removed endpoint heads have vanishing disc
diameter.
\end{enumerate}
All error bounds are uniform over successful outcomes.
\end{proposition}

\begin{proof}
Let the prescribed parameter accuracy and failure
probability decrease to zero in stages. At one stage,
choose \(J\) so that the tail in
\eqref{loc:failure-bound} is smaller than its failure
tolerance. Pass to the subsequence determining the
finite set \(H\). Choose endpoint fences with the
prescribed small parameter size and protect their
access chains. Next choose the \(H\)-caps with
still smaller parameter size and physical cut
diameter \(\rho\) making the first term of
\eqref{loc:failure-bound} small. Only then take
the mesh small enough for all deterministic
contour and cell comparisons at that stage.
Iterate and take a diagonal subsequence.

On success, Lemma~\ref{loc:shield} has removed
all parts beyond \(S+3+K\delta_n\).
The central core and the two front corners
remain, and every new boundary is a wired
extension of one original bank. This proves
the uniform diameter assertion and the
conditional-law assertion. The failure
probability tends to zero by
Lemma~\ref{loc:failure}.

It remains to verify coordinate control
uniformly in the explored outcomes.
In initial disc coordinates, the removed
heads and their filled components are
small attached obstructions by
Lemma~\ref{loc:heads}. Thus their removal
changes normalized coordinates by \(o(1)\)
along the diagonal stages, by
Lemma~\ref{loc:obstructions}.

The portals in the resulting domain also
have uniformly vanishing parameter
diameter. For the \(H\)-portals this follows
from their chosen angular sizes. For
circle-\(S\) portals it follows from
Lemma~\ref{loc:remote}, since the entire
circle is outside \(B(0,S-1)\).
Their noncentral shadows are consequently
small. Every subsequent searched or
discarded portion lies behind one of
these portals. Its points attach to the
old boundary within that small shadow.
Any remaining boundary-cell modifications
have vanishing parameter diameter by
Lemma~\ref{geo:crosscut} and attach
through their incident sectors.
Filling discarded components preserves
the small attachment property.
A second application of
Lemma~\ref{loc:obstructions} therefore
gives uniform \(o(1)\) closeness between
the final and initial normalized maps.

In particular all fixed kernel compacts
are eventually unmodified; the successive
cuts lie in parameter neighborhoods
shrinking to the circle. The two new
marks are within the endpoint fence
neighborhoods, whose diameters tend to
zero. This proves marked Carath\'eodory
convergence of every deterministic
sequence of successful outputs. The
same reasoning is uniform over outcomes,
because only the deterministic stage
tolerances, cell sizes, and parameter
shadows were used.

Finally the original curve on success
is the first head, followed by the
retained interface, followed by the
reverse of the second head. Its
middle part is exactly the ordinary
conditional interface already described.
Uniform coordinate closeness and
vanishing head diameters show that
this concatenation and the retained
curve have distance \(o(1)\) in the
oriented uniform metric modulo
increasing reparametrization. One may
allocate a vanishing parameter interval
to each head, or use continuous
nondecreasing reparametrizations and
approximate them by strictly increasing
ones. This proves the last assertion.
\end{proof}

\begin{corollary}[Universality without a diameter assumption]
\label{loc:universality}
For \(m>0\), every approximation sequence in
Theorem~\ref{thm:main} converges to the law of
Definition~\ref{def:law}, in the specified
normalized-disc curve topology.
\end{corollary}

\begin{proof}
Apply Proposition~\ref{loc:reduction} to any subsequence.
Conditional on a successful record, the retained law
is one of the bounded laws of
Proposition~\ref{uni:bounded}. Their convergence is
uniform over the possible successful outputs in
the required sequential sense. Indeed, if a
bounded continuous test function of curves
had an error bounded below along some sequence
of outputs, choosing one such output at
each mesh would contradict that proposition,
since every chosen domain sequence has
the same marked limit and a common diameter
bound.

Integrate this convergence against the
distribution of the successful records.
The failure probability vanishes, and
the original and retained curves have
vanishing distance on success.
Thus this further subsequence converges
to the same continuum law. Every
subsequence admits such a further
subsequence, which proves full-sequence
convergence. In particular the limiting
law is independent of the lattices,
domain approximations, and nome sequence.
\end{proof}

\paragraph{Negative mass and completion of the theorem.}
For \(m<0\), exchange primal and dual vertices and
the corresponding rhombus diagonals. The prescribed
negative abstract angle is the complement of the
positive abstract angle at the complementary
geometric half-angle. Hence the two sine factors
in the loop weight are exchanged exactly, while
the number of closed loops and the angle bound
are preserved. The two boundary banks exchange
roles. Apply the positive-mass construction on
the dual domain with the starting and terminal
marks interchanged, and reverse the resulting
oriented path. Any boundary half-cell adjustment
has vanishing normalized diameter by
Lemma~\ref{geo:crosscut} and the attachment
argument above. This yields the claimed
continuum characterization and full-sequence
universality for the negative branch as well.
Together with Proposition~\ref{uni:bounded},
Corollary~\ref{loc:universality} proves
Theorem~\ref{thm:main}.
